\printedpage{14}
% head: A. Grothendieck : Dérivateurs, Ch. II, Édition des Universités de Montpellier II et Paris VII
% foot: Version 8 juillet 2001

\[
\begin{tikzcd}
\mathrm{b}' \arrow[r, Rightarrow] & \mathrm{b} \arrow[rr, Rightarrow] \arrow[d, Rightarrow] & & \mathrm{a} \arrow[d, Rightarrow] \\
 & \mathrm{d} \arrow[r, Leftrightarrow] & \mathrm{e} \arrow[r, Leftrightarrow] & \mathrm{c}.
\end{tikzcd}
\]

Let us prove the implications on the bottom line. Let us consider d). If $\begin{tikzcd}[cramped, row sep=0.1em, column sep=0.8em, baseline=(current bounding box.center)] a \arrow[dr] & \\ & x \\ b \arrow[ur] & \end{tikzcd}$ \texttt{[is]} in $S$, to say that the fiber product in $\X$ is ``nonempty'' means, since $S$ is a system of generators and in particular \textit{nonempty,} that there exists $c$ in $S$ and $c\longrightarrow a\times_x b$, or what amounts to the same, $\begin{tikzcd}[cramped, row sep=0.1em, column sep=0.8em, baseline=(current bounding box.center)] & a \\ c \arrow[ur, dotted] \arrow[dr, dotted] & \\ & b \end{tikzcd}$, making the diagram $\begin{tikzcd}[cramped, row sep=0.1em, column sep=0.8em, baseline=(current bounding box.center)] & a \arrow[dr] & \\ c \arrow[ur, dotted] \arrow[dr, dotted] & & x \\ & b \arrow[ur] & \end{tikzcd}$ commutative. Thus d) is equivalent to PS 1, i.e. to e) ($\X$ locally weakly left filtering). Let us consider condition c): the $\X/a=(S/a)^\sim$ are irreducible. This means that for two \textit{nonempty} subobjects $U,V\subseteq a$ in $\X^\wedge$, we have $U\cap V\neq\emptyset$, or else that for two sieves

\msp{page 22}$\underline{U}$, $\underline{V}$ of $S/x$, if $\underline{U}$, $\underline{V}$ nonempty (which means that the subobjects of $x$ in $\X^\wedge$ which they define are ``nonempty'', or that the objects of $S$ are ``nonempty'' in $\X^\wedge$), then $\underline{U}\cap\underline{V}$ is nonempty. We may suppose that $\underline{U}$, $\underline{V}$ is each generated by a single arrow $a\xrightarrow{u}x$, $b\xrightarrow{v}x$. The intersection sieve consists of the objects $c\longrightarrow x$ of $S/x$ such that $c\longrightarrow x$ factors through a $\begin{tikzcd}[cramped, row sep=0.1em, column sep=0.8em, baseline=(current bounding box.center)] & a \\ c \arrow[ur, dotted] \arrow[dr, dotted] & \\ & b \end{tikzcd}$, so to say that this sieve is not empty means that PS 1 is satisfied for the two given arrows $u$, $v$. Thus c) is equivalent to condition PS 1\texttt{[, i.e. to]} e.

Thus we have established
\[
\mathrm{b}'\Longrightarrow\mathrm{b}\Longrightarrow\mathrm{a}\Longrightarrow(\mathrm{c}\Longleftrightarrow\mathrm{d}\Longleftrightarrow\mathrm{e}),
\]
it remains to prove that c $\Longrightarrow$ b$'$. So let $\begin{tikzcd}[cramped, row sep=0.1em, column sep=0.8em, baseline=(current bounding box.center)] U \arrow[dr] & & V \arrow[dl] \\ & W & \end{tikzcd}$ be a diagram of ``nonempty'' objects, with $W$ 0-connected, we must prove that the fiber product is nonempty. Considering $a\longrightarrow U$ and $b\longrightarrow V$ with $a,b\in S$, we are reduced to the case where $U=a$, $V=b$ in $S$. Let us prove that if $a\times_W b$ is ``empty'', then $W$ is disconnected. For this, consider the subcategory $C$ of $S/W$ consisting of the $x\longrightarrow W$ such that $a\times_W x$ is ``nonempty''. Denoting by $w(a)$, $w(x)$ the image of $a$, $x$ in $W$, this means that $w(a)\cap w(x)\neq\emptyset$.

\msp{page 23}Note that any quotient in $\X^\wedge$ of an irreducible object is irreducible, so $w(a)$, $w(x)$ are irreducible objects. This shows that for $x$ in $C$ and $x'$ in $S/W$, we have $x'$ in $C$, i.e. $w(a)\cap w(x')\neq\emptyset$ if and only if $w(x)\cap w(x')\neq\emptyset$. (Direct lemma: in an object $W$ of a topos $\X$, if we have \texttt{[some]} irreducible subobjects $W_1$, $W_2$, $W_3$, if one of the two \texttt{[rather three]} meets the other two, these meet.) This shows that $C$ is a sieve, since $x'\longrightarrow x$ implies $w(x')\subseteq w(x)$, so $w(x')$ meets $w(x)$, hence also $w(a)$. But then the cosieve $C'=(S/W)\backslash C$ is also a sieve, for if $c$
\printedpage{15}
% head: A. Grothendieck : Dérivateurs, Ch. II, Édition des Universités de Montpellier II et Paris VII
% foot: Version 8 juillet 2001
is in $C'$, i.e. \texttt{[if]} $w(c)$ does not meet $w(x)$, a fortiori not any $x'$ over $x$, for which we will have $w(x')\subseteq w(x)$. Thus, $S/W$ is decomposed as a sum of two opens $C$, $C'$, nonempty since one contains $a$, the other $b$. Consequently, $W$ is decomposed as a sum of two nonempty subobjects, so $W$ is not 0-connected.

We have proved 2.8, hence also (returning to 2.6) the implications toward PS 1 in the diagram on page 20. It remains to prove
\[
\left\{
\begin{aligned}
&\mathrm{b}\Longrightarrow\mathrm{PS}\ 2,\ \mathrm{c}'\Longrightarrow\mathrm{PS}\ 2\\
&\mathrm{PS}\ 1+\mathrm{PS}\ 2\ (\text{i.e. d})\Longrightarrow\mathrm{a}'
\end{aligned}
\right.
\]
(\footnote[32]{NB in the end we do not prove c$'\Longrightarrow$ PS 2.}).

\msp{page 24}We therefore place ourselves under the hypothesis where PS 1 is satisfied (since b), c$'$) and d) imply PS 1), thus where we are under the equivalent conditions of 2.8 - $\X$ is locally irreducible. Let us then consider condition PS 2, and prove that it is implied both by b), and by c$'$), and that it implies a$'$). It means that for a double arrow $a\overset{u}{\underset{v}{\rightrightarrows}}b$ in $S$, the cokernel \texttt{[rather kernel]} in $S^\wedge$, or, what amounts to the same, in $\X^\wedge=S^\sim\subseteq S^\wedge$, is ``nonempty''. We can write it in the form of a fiber product
\[
\begin{tikzcd}
a \arrow[r, "{\scriptstyle (u,v)}"] & b\times b \\
\operatorname{Ker}(u,v) \arrow[u] \arrow[r] & b. \arrow[u, "\mathrm{diag}"']
\end{tikzcd}
\]
Here by hypothesis $a$ and $b$ are 0-connected objects of $\X^\wedge$ (those of $S$ being so), c) holds (($\X$ locally irreducible and) its connected components \texttt{[are]} totally 0-connected), then $b\times b$ is 0-connected as a product of two 0-connected objects over the \textit{same connected component,} which is totally 0-connected. Since $\X\times\X$ is locally irreducible, it follows that the fiber product is nonempty (2.8), hence b) $\Longrightarrow$ d). Let us prove d) $\Longrightarrow$ a), i.e. if

\msp{page 25}$S$ is locally cofiltering, then the induced topoi $\X/W$, where $W$ is 0-connected, are totally 0-connected. This induced topos is described by the site $S/W$, i.e. $S/W$ is a system of generators. Now if $S$ is locally filtering, $S/W$ (étale over $S$) is so too. So we may suppose $W=e$, and to prove that if $S$ is locally cofiltering and $\X$ 0-connected, then $\X$ is totally 0-connected. Note that $\X$ 0-connected is equivalent to $S$ 0-connected - one should make this explicit all the same:

\textbf{Lemma 2.9.} \textit{Let $\X$ be a topos, $S\subseteq\X^\wedge$ a full generating subcategory consisting of objects} that are 0-connected. \textit{Then $\X$ is $0$-connected if and only if $S$ is - more generally, the ordered set of both open and closed subobjects of $S$ in} Cat \textit{(or else of the
\printedpage{16}
% head: A. Grothendieck : Dérivateurs, Ch. II, Édition des Universités de Montpellier II et Paris VII
% foot: Version 8 juillet 2001
open subobjects $=$ sieves of $S$ that are also closed, i.e. cosieves) is isomorphic to that of the direct summands of $\X$.}

I used it tacitly in the proof of 2.8, for the induced topos $\X/W$. I do not spell out the proof.

Thus (returning to the proof of d) $\Longrightarrow$ c) in 2.6.) $S$ being 0-connected and locally cofiltering,

\msp{page 26}\texttt{[it]} is in fact cofiltering. So for two objects $a$, $b$ of $S$, $\exists\ c$ with $c\longrightarrow a$, $c\longrightarrow b$ (condition PS 1$'$, stronger than PS 1). To prove that $\X$ is totally 0-connected, we are reduced to proving by 2.3 that the $a\times b$ in $\X^\wedge$ are 0-connected. By 2.9, it amounts to the same to prove that they are 0-connected in $S^\wedge$. So we are reduced to proving the

\textbf{Lemma 2.10.} \textit{Let $S$ be a cofiltering category. If $a$, $b$ are two objects of $S$, $a\times b$ is an object that is} $0$\textit{-connected in $S^\wedge$.}

We see at once that it suffices to prove (and in fact it amounts to the same, taking into account that $a\times b$ is locally irreducible) that if we have two arrows
\[
\begin{tikzcd}
x \arrow[r] & a\times b \\
 & y, \arrow[u]
\end{tikzcd}
\]
then $x\times_{a\times b}y\neq\emptyset$, i.e. $\exists\ z$ in $S$ making commutative
\[
\begin{tikzcd}
x \arrow[r] & a\times b \\
z \arrow[u] \arrow[r] & y. \arrow[u]
\end{tikzcd}
\]
Now by PS 1$'$ there does indeed exist a $z$ mapping to $x$ and to $y$, but the two composite arrows
\[
z\rightrightarrows a\times b,
\]
i.e. each of the double arrows
\[
z\rightrightarrows a \tag*{$(*)$}
\]
\[
z\rightrightarrows b \tag*{$(*')$}
\]
may not be equal. But by PS 2,

\msp{page 27}we can find
\[
z'\longrightarrow z
\]
\printedpage{17}
% head: A. Grothendieck : Dérivateurs, Ch. II, Édition des Universités de Montpellier II et Paris VII
% foot: Version 8 juillet 2001
equalizing the double arrow $(*)$, so we may suppose that this one is already trivial. Applying PS 2 once again, by a $z''\longrightarrow z'$, we similarly equalize the double arrow $(*')$, and we win. We have therefore proved, finally,
\[
\begin{tikzcd}
\mathrm{a} \arrow[r, Leftrightarrow] \arrow[d, Rightarrow] & \mathrm{a}' \arrow[r, Rightarrow] & \mathrm{b} \arrow[r, Leftrightarrow] & \mathrm{d}\ (=\mathrm{PS}\ 1+\mathrm{PS}\ 2) \\
\mathrm{c} \arrow[r, Leftrightarrow] & \mathrm{c}',
\end{tikzcd}
\]
but not the reverse implication. In other words, if the induced topoi $\X/a$, for $a$ in $S$, are totally 0-connected, I have not proved that the same holds for the $\X/W$, $W$ a 0-connected object \textit{whatsoever} of $\X^\wedge$. Thus, we are at the following point

\textbf{Perplexity 2.11.} \textit{Let $X$ be in} Cat. \textit{If $X$ is locally left filtering, then the fiber products $a\times_x b$ of objects in $X$ are objects that are} $0$\textit{-connected in $X^\wedge$, and this condition in turn implies PS 1, i.e. that $X$ is locally irreducible. But in the presence of this hypothesis PS 1, or $X$ locally irreducible, the condition on the fiber products $a\times_x b$ is equivalent to the following: for every diagram} \texttt{[in]} \textit{solid lines}

\msp{page 28}
\[
\begin{tikzcd}
c \arrow[r] & a\times_x b \\
e \arrow[u, dotted] \arrow[r, dotted] & d, \arrow[u]
\end{tikzcd}
\qquad\qquad
\textit{with $x$, $a$, $b$, $c$, $d$ in $X$,}
\]
\textit{we can complete it to a diagram that is} commutative, \textit{with $e$ in $X$. In other words, every commutative diagram in solid lines below can be completed by adding to it an ``initial object'' $e$} % [sic] d'autre termes -> d'autres termes, missing plural s
\[
\begin{tikzcd}[column sep=small, row sep=small]
 & & a \arrow[ddr] & \\
 & c \arrow[ur] \arrow[dddr] & & \\
e \arrow[ur, dotted] \arrow[dr, dotted] & & & x \\
 & d \arrow[uuur] \arrow[dr] & & \\
 & & b \arrow[uur] &
\end{tikzcd}
\]
(\footnote[33]{It suffices to require for $e$ the commutativity of the two squares $(e\ c\ b\ d)$ and $(e\ d\ a\ c)$.}).

This condition is thus implied by $X$ locally cofiltering, and from what precedes it is also satisfied more generally if and only if (\footnote[34]{if and only if?}) the $X/x$ are cofiltering (\footnote[35]{It means that $X$ satisfies PS 1, and the condition

\textbf{PS$'$ 2} Every commutative diagram
\[
a\rightrightarrows b\longrightarrow x
\]
% footnote continued on printed page 18
fits into a commutative diagram
\[
y\longrightarrow a\rightrightarrows b\longrightarrow x.
\]
}). But is the converse true - does this condition of total 0-connectedness of the $X/x$ imply that $X$ is totally 0-connected? But I have not constructed an example where the $X/x$
\printedpage{18}
% head: A. Grothendieck : Dérivateurs, Ch. II, Édition des Universités de Montpellier II et Paris VII
% foot: Version 8 juillet 2001
are cofiltering without $X$ being locally cofiltering (or what amounts to the same, of $X$ 0-connected with the $X/x$ cofiltering without $X$ thereby being cofiltering)(\footnote[36]{But there is an obvious \textit{counterexample}, if $X$ is a 0-connected groupoid, the $X/x$ are $\approx e$, hence cofiltering, but $X$ is cofiltering (i.e. \texttt{[illegible word(s)]}) only if $X$ is 1-connected. E.g. if $C$ is a groupoid $\neq 1$, then $X=B_G$ is such that the $X/x$ are cofiltering, but $X$ not cofiltering. % [sic] si $C$ est un groupoïde -> si $G$ est un groupoïde, the groupoid is then written $B_G$
}) ; nor an example where the $X/x$ are totally 0-connected, without the $X/x$ being cofiltering. But let us note the

\textbf{Corollary 2.12:} \textit{The equivalent conditions a), a$'$), b), d) (excluding c), c$'$), apparently weaker) of theorem 2.6}

\msp{page 29}\textit{are also equivalent to the following}
\begin{itemize}[leftmargin=2.5em]
\item[\textit{e)}] \textit{$\forall\ x$ in $S$, the induced topos $\X/x$ and the induced topos $\X/x\times x$ is irreducible} (\footnote[37]{NB $\X/x$ unnecessary.}),
\end{itemize}
\textit{or else to the following} % [sic] où encore -> ou encore, conjunction "ou" without accent
\begin{itemize}[leftmargin=2.5em]
\item[\textit{e$'$)}] \textit{$\forall\ x$ in $S$, $S^\wedge/x$ and $S^\wedge/x{\times}x$ are irreducible, i.e. $S/x$ and $S/x{\times}x$ are irreducible} (\footnote[38]{$S/x$ unnecessary, for $x\xrightarrow{\mathrm{diag}}x\times x$ is a nonempty open immersion, so if $x\times x$ is irreducible, $x$ is too. But it amounts to the same to say that the $S/x$ are irreducible (condition PS$^*$ 1, cf. 2.8, p. 21) \textit{and} that the $x\xrightarrow{\mathrm{diag}}x\times x$ are dominant (which indeed implies that $x\times x$ is irreducible). This immediately gives back the form PS$^*$ 2 of the cofiltering condition. % [sic] ce que implique -> ce qui implique, relative pronoun subject
}).
\end{itemize}
This is the apparently weakest ``topological'' form that I can find for the condition $S$ locally cofiltering, i.e. $\X$ (or Top$(S)$, at will) locally totally 0-connected.

\textbf{Definition 2.13.} We say that a topos $\X$ is \textit{locally totally} $0$\textit{-connected} if it satisfies the equivalent conditions a, a$'$, b, d of theorem 2.6, i.e. if $\X$ is locally connected and if in addition every 0-connected induced topos is totally 0-connected - the fiber product of two 0-connected objects over a third is 0-connected. It amounts to the same to say that every (or that one) full generating subcategory $S$ of $\X^\wedge$ consisting of 0-connected objects is locally cofiltering (still assuming $\X$ locally connected, or what amounts to the same, the existence of an $S$ consisting of 0-connected objects).

We have not yet finished the proof of 2.6 - it remains to prove that if $\X$ (assumed locally totally 0-connected) is 0-connected, it is $\infty$-connected - or else that the connected components of $\X$, assumed locally totally 0-connected, are $\infty$-connected, i.e. $W_\infty$-aspheric.

\msp{page 30}\textbf{Corollary 2.14.} \textit{Let us note all the same that if $\X$ is of the form Top$(T)$, $T$ a topological space, then $\X$ is locally totally} $0$\textit{-connected if and only if it is locally irreducible, i.e. if and only if $T$ is locally irreducible.}